\section{Cut scan on the 10\% skim}
\label{sec:scan}

Sixteen selections (S0--S15) combine the building blocks of Sec.~\ref{sec:method}, ordered from loose to tight
(Table~\ref{tab:scan}, Figs.~\ref{fig:scan1} and~\ref{fig:scan2}). Each is fitted with the fixed 15\,MeV Gaussian at the
\Dzero{} mass, the free fit and the sideband cross-check.

\begin{table}[H]
\centering\footnotesize
\setlength{\tabcolsep}{4pt}
\caption{Hand cut scan, 10\% skim. $N$: candidates; $S$, $z$: fixed fit (15\,MeV Gaussian at $m_{\mathrm{D}^0}$, pol2,
amplitude unbounded); $\chi^2/$ndf of that fit; $S_\mathrm{sb}$: sideband cross-check; free fit: mean (GeV/$c^2$) /
width (MeV), $\dagger$ = a parameter at its limit (not interpreted). Statistical errors only. Source:
\fpath{logs/scan.txt}.}
\label{tab:scan}
\begin{tabular}{llrrrrrl}
\toprule
 & selection & $N$ & $S$ & $z$ & $\chi^2/$ndf & $S_\mathrm{sb}$ & free fit \\
\midrule
\input{tables/scan_results.tex}
\bottomrule
\end{tabular}
\end{table}

\paragraph{Result.} The pol2 background describes the loosest selections badly ($\chi^2/\mathrm{ndf}\ge2$ for
\scanBadSels; up to \scanChiMax), and their apparent excesses (S0: $z=\scanZmax$) follow the background curvature
(Sec.~\ref{sec:trial}). Among the \scanGoodN{} selections with $\chi^2/\mathrm{ndf}<2$ the largest significance is
$z=\scanZmaxGood$ (\scanZmaxGoodSel), and the tight selections (C, strict PID, $\pt>2$) give $z$ between \scanZminAll{}
and about zero. The free fits end at a parameter limit in \scanFreeLim{} of \scanN{} selections; where they converge
inside the limits (S5, S6, S14, S15) their means are 1.872--1.880\,GeV/$c^2$ (Table~\ref{tab:scan}), above the PDG mass,
with yields compatible with the fixed fit within its error (read from the table). No selection shows a narrow peak.

\begin{figure}[H]
\centering
\includegraphics[width=\textwidth]{scan_mass_1}
\caption{$K\pi$ mass spectra of selections S0--S7 (10\% skim, 5\,MeV bins, statistical errors). Red: fixed fit
(15\,MeV Gaussian at the \Dzero{} mass, dotted line, plus pol2); the fitted yield is printed in each panel.}
\label{fig:scan1}
\end{figure}

\begin{figure}[H]
\centering
\includegraphics[width=\textwidth]{scan_mass_2}
\caption{As Fig.~\ref{fig:scan1} for S8--S15.}
\label{fig:scan2}
\end{figure}
